\documentclass[11pt,a4paper]{article}

\usepackage[T1]{fontenc}
\usepackage[utf8]{inputenc}
\usepackage{lmodern}
\usepackage[margin=2.4cm]{geometry}
\usepackage{amsmath,amssymb,bm}
\usepackage{booktabs}
\usepackage{xcolor}
\usepackage{listings}
\usepackage{mdframed}
\usepackage[colorlinks=true,linkcolor=NavyBlue,citecolor=NavyBlue,urlcolor=NavyBlue]{hyperref}
\usepackage[dvipsnames]{xcolor}
\usepackage{microtype}

\definecolor{codebg}{HTML}{F4F4F0}
\definecolor{coderule}{HTML}{C9C7BC}
\definecolor{keyc}{HTML}{7B3F00}
\definecolor{comc}{HTML}{6B7A6B}

\lstdefinestyle{f90}{
  language=[90]Fortran,
  basicstyle=\ttfamily\scriptsize,
  keywordstyle=\color{keyc},
  commentstyle=\color{comc}\itshape,
  columns=fullflexible,
  keepspaces=true,
  breaklines=true,
  breakatwhitespace=false,
  postbreak=\mbox{\textcolor{coderule}{$\hookrightarrow$}\space},
  showstringspaces=false,
  frame=none,
  backgroundcolor=\color{codebg},
  xleftmargin=4pt,
  framexleftmargin=4pt,
  framexrightmargin=4pt,
  framextopmargin=3pt,
  framexbottommargin=3pt,
}
\lstset{style=f90}

% --- a "code line" block: source reference, the Fortran, then the maths -----------------
\newenvironment{codeline}[1]
  {\par\medskip\noindent\textcolor{coderule}{\rule{\linewidth}{0.4pt}}\par
   \nopagebreak\noindent{\ttfamily\footnotesize\textcolor{keyc}{#1}}\par\nopagebreak\smallskip}
  {\par\medskip}

\newcommand{\code}[1]{\texttt{\small #1}}
\newcommand{\bhat}{\hat{\bm{b}}}
\newcommand{\nhat}{\hat{\bm{n}}}
\newcommand{\that}{\hat{\bm{t}}}
\newcommand{\vE}{\bm{v}_E}
\newcommand{\eph}{\bm{e}_\varphi}

\newmdenv[
  linecolor=coderule,linewidth=0.6pt,
  backgroundcolor=codebg,
  innertopmargin=6pt,innerbottommargin=6pt,
  innerleftmargin=8pt,innerrightmargin=8pt,
  skipabove=8pt,skipbelow=8pt,
]{keybox}

\title{\vspace{-1.5cm}The parallel-velocity boundary condition in JOREK,\\
       with the $\bm{E}\times\bm{B}$ drift term written out}
\author{Máté Szűcs}
\date{27 August 2026}

\begin{document}
\maketitle

\begin{abstract}
\noindent
This note writes out, line by line, the Mach-1 (Bohm) boundary condition on the parallel
velocity as it is actually implemented in JOREK's \code{model600}, including the
$\bm{E}\times\bm{B}$ drift term that JOREK already carries.  Every expression is paired with
the exact source line it comes from, in
\code{models/model600/mod\_boundary\_conditions.f90} at commit \code{9508effe0} of branch
\code{thermoelectric-ohm}.  The last two sections compare the result with the
drift-compatible Bohm condition used by SOLPS-ITER and report a measurement, taken with the
\code{diag\_mach1} A/B diagnostic, showing that the two differ by exactly one factor of
$|\bm{B}|$.
\end{abstract}

\section{The velocity representation}

\code{model600} is a reduced-MHD model.  The velocity is split into a poloidal
$\bm{E}\times\bm{B}$ part written through a stream function $u$, and a parallel part written
through a variable $v_{\rm par}$:
\begin{equation}
  \bm{v} \;=\; \underbrace{R\,\nabla u \times \eph}_{\textstyle \vE}
             \;+\; v_{\rm par}\,\bm{B} .
  \label{eq:ansatz}
\end{equation}

Two things about this ansatz drive everything below.

\begin{keybox}
\textbf{(i)} The stream-function convention fixes the sign of the potential:
$\Phi = -F_0 u$ in \code{model600}.  (The JOREK reference papers write
$\bm{v}_{\rm pol} = -R\nabla u\times\eph$; the code uses $+$.)

\textbf{(ii)} $v_{\rm par}$ is \emph{not} the parallel velocity.  The physical parallel
velocity is
\[
   v_\parallel \;=\; v_{\rm par}\,|\bm{B}| \;=\; v_{\rm par}\,B_{\rm tot}.
\]
Every factor of $1/B_{\rm tot}$ in the code exists to undo this.  Losing track of it is the
single easiest mistake to make in this routine, and \S\ref{sec:compare} is about exactly
that.
\end{keybox}

Writing \eqref{eq:ansatz} in components, with $\bm{x}=(R,Z)$ poloidal:
\begin{equation}
  \vE \;=\; R\,\bigl(-\partial_Z u,\; \partial_R u\bigr),
  \qquad
  \bm{B}_{\rm pol} \;=\; \frac{1}{R}\bigl(\partial_Z\psi,\;-\partial_R\psi\bigr),
  \qquad
  \bm{B}_\varphi = \frac{F_0}{R}\,\eph .
  \label{eq:components}
\end{equation}
Note the structural symmetry: $\vE$ is $-R^2$ times the same differential operator applied to
$u$ that produces $\bm{B}_{\rm pol}$ from $\psi$.  That symmetry is what makes the identity in
\S\ref{sec:identity} come out as cleanly as it does.

The $\vE$ components appear literally in the diagnostic block, which is the easiest place to
read the sign convention off the source:

\begin{codeline}{mod\_boundary\_conditions.f90:637--641}
\begin{lstlisting}
u0_x_d = (   Z_t * u0_s_d - Z_s * u0_t_d ) / xjac      ! = du/dR
u0_y_d = ( - R_t * u0_s_d + R_s * u0_t_d ) / xjac      ! = du/dZ
vE_R_d = - BigR * u0_y_d
vE_Z_d = + BigR * u0_x_d
\end{lstlisting}
\end{codeline}

\section{Geometry at the boundary}

The boundary of a bicubic element is a curve of constant $s$ or constant $t$.  The routine
sets a single index \code{iv\_dir} that selects the derivative \emph{along} that curve:

\begin{codeline}{mod\_boundary\_conditions.f90:254, 259}
\begin{lstlisting}
        iv_dir      = 2      ! t_constant_boundary: differentiate w.r.t. s
        iv_dir      = 3      ! s_constant_boundary: differentiate w.r.t. t
\end{lstlisting}
\end{codeline}

\begin{keybox}
\code{iv\_dir} is the \textbf{tangential} derivative index, not the normal one.  On a
\code{t\_constant} boundary $s$ runs along the wall, so \code{iv\_dir = 2} means
$\partial/\partial s$, which is tangential.  (\code{iv\_perp\_dir}, used elsewhere in the same
file, is the normal one.)  This is worth stating explicitly because reading it the other way
inverts the conclusion of \S\ref{sec:identity}.
\end{keybox}

Throughout, a subscript $b$ denotes this tangential derivative, scaled to the element by
\code{element\_size\_0}.  So for any field $f$,
\begin{equation}
  f_{0b} \;\equiv\; \code{f0\_b} \;=\; \partial_b f \;=\; \nabla f \cdot \that \;\; \Delta\ell ,
\end{equation}
where $\that$ is the unit tangent and $\Delta\ell = \code{dl} = \sqrt{R_b^2+Z_b^2}$ is the
arc length per unit of the local boundary coordinate.

\begin{codeline}{mod\_boundary\_conditions.f90:468, 590, 583}
\begin{lstlisting}
ps0_b   = node_list%node(inode)%values(1,iv_dir,var_psi)  * element_size_0
u0_b    = node_list%node(inode)%values(1,iv_dir,var_u)    * element_size_0
Vpar0_b = node_list%node(inode)%values(1,iv_dir,var_Vpar) * element_size_0
\end{lstlisting}
\end{codeline}

\noindent The remaining geometric quantities:

\begin{codeline}{mod\_boundary\_conditions.f90:512--521}
\begin{lstlisting}
normal     = normal / norm2(normal)                                   ! outward pointing
direction  = sign(1.d0,dot_product((/ps0_y,-ps0_x/),normal))
Btot       = sqrt(F0**2 + ps0_x**2 + ps0_y**2) / BigR
dl         = sqrt(R_b**2 + Z_b**2)
bn         = dot_product( (/ps0_y,-ps0_x/), normal ) /  (BigR*Btot)   ! B.n/Btot
\end{lstlisting}
\end{codeline}

\begin{align}
  B_{\rm tot} &= \frac{\sqrt{F_0^2 + |\nabla\psi|^2}}{R} = |\bm{B}| ,
  &
  b_n &= \frac{\bm{B}\cdot\nhat}{|\bm{B}|} = \frac{\bm{B}_{\rm pol}\cdot\nhat}{B_{\rm tot}} ,
  &
  \sigma &\equiv \code{direction} = \operatorname{sign}(b_n) .
  \label{eq:geom}
\end{align}
The toroidal field has no normal component at the wall, so only $\bm{B}_{\rm pol}$ enters
$b_n$.  Comparing the \code{bn} line with \eqref{eq:components} gives the sign relation used
below,
\begin{equation}
  \bm{B}_{\rm pol}\cdot\nhat \;=\; b_n\,B_{\rm tot} .
  \label{eq:bpoln}
\end{equation}

\subsection*{The obliqueness factor}

\code{factor} is a smoothing weight, active only where the normal field component changes sign
between the two ends of a boundary edge ($\code{bn\_1}\cdot\code{bn\_2} < 0$), i.e.\ at a
tangency point.  Away from such points it is identically $1$.

\begin{codeline}{mod\_boundary\_conditions.f90:534--554}
\begin{lstlisting}
c_1 = vpar_smoothing_coef(1); c_2 = vpar_smoothing_coef(2); c_3 = vpar_smoothing_coef(3)
if ((vpar_smoothing) .and. (bn_1*bn_2 .lt. 0.d0)) then
  if (c_2 .gt. 0d0) then
    factor    = 0.25d0 * ( 1.d0 + tanh( (abs(bn) - c_1) / c_2 ) )**2 - c_3
  else
    factor    = tanh(bn/c_1)
    direction = 1.d0
  endif
else
  factor    = 1.d0
endif
Hfact_b   = factor * R_b / BigR  + factor_b
\end{lstlisting}
\end{codeline}

\begin{equation}
  \mathcal{F} \;\equiv\; \code{factor} \;=\;
  \begin{cases}
    \dfrac{1}{4}\Bigl[1 + \tanh\dfrac{|b_n| - c_1}{c_2}\Bigr]^{2} - c_3, & c_2 > 0,\\[2.2ex]
    \tanh\!\bigl(b_n/c_1\bigr), \quad \sigma \to 1, & c_2 \le 0,\\[1.2ex]
    1 & \text{away from tangency.}
  \end{cases}
  \label{eq:factor}
\end{equation}

\code{Hfact\_b} is not a separate physical quantity: it is the logarithmic derivative of the
whole prefactor $\mathcal{F}/B_{\rm tot}$.  Under the routine's stated approximation that the
field is frozen (so $B_{\rm tot}\simeq F_0/R$ and $\partial_b B_{\rm tot}^{-1} =
R_b/F_0 = R_b/(R\,B_{\rm tot})$),
\begin{equation}
  \partial_b\!\left(\frac{\mathcal{F}}{B_{\rm tot}}\right)
  \;=\; \frac{1}{B_{\rm tot}}\underbrace{\left(\mathcal{F}\frac{R_b}{R} + \mathcal{F}_b\right)}
        _{\textstyle \code{Hfact\_b}} .
  \label{eq:hfact}
\end{equation}

\section{The boundary condition itself}

\subsection{The value row}

\begin{codeline}{mod\_boundary\_conditions.f90:617 \quad\normalfont\itshape(the line this note is about)}
\begin{lstlisting}
Mach1BC = - Vpar0 + direction / Btot * factor * cs0 + factor / Btot * BigR**2 * U0_b/ps0_b
\end{lstlisting}
\end{codeline}

\begin{equation}
  \boxed{\;
  \mathcal{M} \;=\; -\,v_{\rm par}
  \;+\; \frac{\sigma\,\mathcal{F}}{B_{\rm tot}}\,c_s
  \;+\; \frac{\mathcal{F}}{B_{\rm tot}}\,R^{2}\,\frac{\partial_b u}{\partial_b \psi}
  \;=\; 0
  \;}
  \label{eq:mach1}
\end{equation}

with the sound speed and its temperature derivatives

\begin{codeline}{mod\_boundary\_conditions.f90:612--615}
\begin{lstlisting}
cs0      =   sqrt(gamma*(Ti0+Te0))
cs0_T    =   0.5d0  * gamma    / cs0
cs0_TT   = - 0.25d0 * gamma**2 / cs0**3
cs0_TTT  = 3.d0/8.d0* gamma**3 / cs0**5
\end{lstlisting}
\end{codeline}

\begin{equation}
  c_s = \sqrt{\gamma\,(T_i + T_e)},\qquad
  \frac{\partial c_s}{\partial T} = \frac{\gamma}{2c_s},\qquad
  \frac{\partial^2 c_s}{\partial T^2} = -\frac{\gamma^2}{4c_s^3},\qquad
  \frac{\partial^3 c_s}{\partial T^3} = \frac{3\gamma^3}{8c_s^5}.
\end{equation}

The third term of \eqref{eq:mach1} is the drift term.  It is the subject of the next section.

\subsection{The tangential-derivative row}

The trace space of a bicubic $C^1$ element carries a value \emph{and} a tangential derivative
at each node, so the condition must be imposed on both.  Differentiating \eqref{eq:mach1}
along the boundary:

\begin{codeline}{mod\_boundary\_conditions.f90:621--622, 655}
\begin{lstlisting}
dMach1BC = - Vpar0_b + direction / Btot * factor  * cs0_T * (Ti0_b+Te0_b)  &
                     + direction / Btot * Hfact_b * cs0
! --- only for n_order >= 5:
dMach1BC = dMach1BC + factor / Btot * BigR**2 * U0_bb/ps0_b
\end{lstlisting}
\end{codeline}

\begin{equation}
  \partial_b\mathcal{M} \;=\; -\,\partial_b v_{\rm par}
  \;+\; \frac{\sigma\mathcal{F}}{B_{\rm tot}}\frac{\partial c_s}{\partial T}\bigl(\partial_b T_i + \partial_b T_e\bigr)
  \;+\; \frac{\sigma}{B_{\rm tot}}\,\code{Hfact\_b}\;c_s
  \;+\; \underbrace{\frac{\mathcal{F}}{B_{\rm tot}}R^2\frac{\partial_b^2 u}{\partial_b\psi}}
        _{\textstyle n_{\rm order}\ge 5\ \text{only}} .
  \label{eq:dmach1}
\end{equation}

\begin{keybox}
\textbf{Approximations in \eqref{eq:dmach1}, stated by the code itself} (line 556: \emph{``For
the BC's the magnetic field is assumed constant, and we do not take derivatives of psi into
account''}):
\begin{itemize}\itemsep2pt
  \item $\partial_b\psi$, $\partial_b R^2$ and $\partial_b B_{\rm tot}$ are not differentiated
        in the drift term.  Only $R_b/R$ survives, through \code{Hfact\_b}, and only in the
        $c_s$ term.
  \item At $n_{\rm order} = 3$ --- the production setting --- \textbf{the drift term does not
        contribute to the derivative row at all}, because $\partial_b^2 u$ is not a degree of
        freedom.  The Bohm condition is drift-corrected in its value but not in its
        tangential slope.
\end{itemize}
\end{keybox}

\subsection{The second-derivative row ($n_{\rm order}\ge5$)}

\begin{codeline}{mod\_boundary\_conditions.f90:657--668}
\begin{lstlisting}
d2Mach1BC    = - Vpar0_bb + direction / Btot * factor   * cs0_TT * (Ti0_b+Te0_b)**2   &
                          + direction / Btot * factor   * cs0_T  * (Ti0_bb+Te0_bb)   !&
                          !+ direction / Btot * Hfact_b  * cs0_T  * T0_b *2.0 !&
                          !+ direction / Btot * Hfact_bb * cs0
\end{lstlisting}
\end{codeline}

\begin{equation}
  \partial_b^2\mathcal{M} \;=\; -\,\partial_b^2 v_{\rm par}
  \;+\; \frac{\sigma\mathcal{F}}{B_{\rm tot}}\left[
        \frac{\partial^2 c_s}{\partial T^2}\bigl(\partial_b T_i + \partial_b T_e\bigr)^2
      + \frac{\partial c_s}{\partial T}\bigl(\partial_b^2 T_i + \partial_b^2 T_e\bigr)
        \right].
\end{equation}
The two \code{Hfact} contributions are commented out in the source, and the drift term has no
second-derivative counterpart at all.  This row is therefore the least complete of the three;
it is unused at $n_{\rm order}=3$.

\subsection{Jacobian columns}

The condition is imposed by row replacement with a large weight \code{zbig}.  The linearisation
treats $v_{\rm par}$, $T$ and $u$ as the only variables --- $\psi$, and with it the whole
geometry, is frozen inside a Newton iteration.

\begin{codeline}{mod\_boundary\_conditions.f90:618--620, 623--630, 656}
\begin{lstlisting}
Mach1BC_v   = - 1.0
Mach1BC_T   =           + direction / Btot * factor  * cs0_T
Mach1BC_u   =                                        + factor / Btot * BigR**2 * element_size_0/ps0_b
dMach1BC_v  = - element_size_0
dMach1BC_T  =           + direction / Btot * factor  * cs0_TT* T0_b  &
                        + direction / Btot * Hfact_b * cs0_T
dMach1BC_Tb =           + direction / Btot * factor  * cs0_T * element_size_0
dMach1BC_ubb=           + factor / Btot * BigR**2 * element_size_3/ps0_b
\end{lstlisting}
\end{codeline}

\begin{center}
\begin{tabular}{llll}
\toprule
row & column & code & expression \\
\midrule
value & $v_{\rm par}$ & \code{Mach1BC\_v}    & $-1$ \\
value & $T$           & \code{Mach1BC\_T}    & $\sigma\mathcal{F}B_{\rm tot}^{-1}\,\partial_T c_s$ \\
value & $u$           & \code{Mach1BC\_u}    & $\mathcal{F}B_{\rm tot}^{-1}R^2\,h_0/\partial_b\psi$ \\
$\partial_b$ & $v_{\rm par}$ & \code{dMach1BC\_v}  & $-h_0$ \\
$\partial_b$ & $T$           & \code{dMach1BC\_T}  & $\sigma B_{\rm tot}^{-1}\bigl[\mathcal{F}\,\partial^2_T c_s\,\partial_b T + \code{Hfact\_b}\,\partial_T c_s\bigr]$ \\
$\partial_b$ & $\partial_b T$& \code{dMach1BC\_Tb} & $\sigma\mathcal{F}B_{\rm tot}^{-1}\,\partial_T c_s\, h_0$ \\
\bottomrule
\end{tabular}
\end{center}

\noindent with $h_0 = \code{element\_size\_0}$.  Note that \code{Mach1BC\_u} is the drift
term's Jacobian entry --- it already exists, which matters for \S\ref{sec:fix}.

\section{The identity: what the drift term actually is}
\label{sec:identity}

The drift term in \eqref{eq:mach1} is written in element coordinates,
$R^2\,\partial_b u/\partial_b\psi$, which hides its meaning.  Converting it costs four lines.

Take the unit tangent oriented as $\that = (n_Z, -n_R)$.  From \eqref{eq:components},
\begin{align}
  \vE\cdot\nhat
    &= R\bigl(-u_Z n_R + u_R n_Z\bigr)
     = R\,\nabla u\cdot\that
     = R\,\frac{\partial_b u}{\Delta\ell},
     \label{eq:vEn}\\[4pt]
  \bm{B}_{\rm pol}\cdot\nhat
    &= \frac{1}{R}\bigl(\psi_Z n_R - \psi_R n_Z\bigr)
     = -\frac{1}{R}\,\nabla\psi\cdot\that
     = -\frac{1}{R}\,\frac{\partial_b\psi}{\Delta\ell}.
     \label{eq:Bn}
\end{align}
The arc-length scale $\Delta\ell$ is common to both.  Dividing and using \eqref{eq:bpoln}:

\begin{keybox}
\begin{equation}
  R^{2}\,\frac{\partial_b u}{\partial_b \psi}
  \;=\; -\,\frac{\vE\cdot\nhat}{\bm{B}_{\rm pol}\cdot\nhat}
  \;=\; -\,\frac{\vE\cdot\nhat}{b_n\,B_{\rm tot}} .
  \label{eq:identity}
\end{equation}
\end{keybox}

\noindent This is why \eqref{eq:mach1} works at all: the ratio of two tangential derivatives
of two different fields is a \emph{normal} quantity, because differentiating along the
boundary and dotting the corresponding vector into the normal are the same operation up to the
same scale factor for both $u$ and $\psi$.  It also explains why the term is written this way:
$\partial_b u$ and $\partial_b\psi$ are both directly available as nodal degrees of freedom on
the boundary, whereas $\vE\cdot\nhat$ and $b_n$ are not.

Substituting \eqref{eq:identity} into \eqref{eq:mach1} and multiplying by $B_{\rm tot}$ to
convert $v_{\rm par}\to v_\parallel$:

\begin{keybox}
\begin{equation}
  \boxed{\;
  v_\parallel^{\,\rm JOREK}
  \;=\; \mathcal{F}\left[\,\sigma\,c_s \;-\; \frac{\vE\cdot\nhat}{b_n\,B_{\rm tot}}\,\right]
  \;}
  \label{eq:jorekvpar}
\end{equation}
\end{keybox}

\noindent So JOREK's Mach-1 condition is \emph{already} drift-aware: the second term has the
$\vE\cdot\nhat/b_n$ structure of a drift correction, not merely some numerical smoothing.

\section{Comparison with the drift-compatible (SOLPS) condition}
\label{sec:compare}

The Bohm condition constrains the \emph{total} flux into the wall, not the parallel flux.  With
drifts retained, the total normal velocity is
\begin{equation}
  \bm{v}\cdot\nhat \;=\; v_\parallel\,b_n \;+\; \vE\cdot\nhat \;\;\ge\;\; \sigma\,c_s\,b_n ,
\end{equation}
so the condition to impose on $v_\parallel$ is
\begin{equation}
  \boxed{\;
  v_\parallel^{\,\rm req}
  \;=\; \sigma\,c_s \;-\; \frac{\vE\cdot\nhat}{b_n}
  \;}
  \label{eq:required}
\end{equation}
which is the form SOLPS-ITER imposes (manual 3.0.9, \code{BCMOM} = 3/13), together with a clip
at $\pm 2c_s|b_x|$ to keep the correction bounded near tangency.

Comparing \eqref{eq:jorekvpar} with \eqref{eq:required}, the $\sigma c_s$ terms agree (up to
$\mathcal{F}$, which is $1$ away from tangency), and the drift terms differ by exactly one
factor $B_{\rm tot}$:
\begin{equation}
  \frac{\text{JOREK drift term}}{\text{required drift term}}
  \;=\; \frac{\mathcal{F}}{B_{\rm tot}} .
  \label{eq:ratio}
\end{equation}

\begin{keybox}
\textbf{Reading.} Both terms in \eqref{eq:mach1} carry the same $1/B_{\rm tot}$ prefactor.
That is \emph{correct} for the $c_s$ term, because $v_\parallel = v_{\rm par}B_{\rm tot}$
converts it into $\mathcal{F}\sigma c_s$.  The drift term needed the $B_{\rm tot}$ to survive
the same conversion, and it does not have it.  This reads as a missing $B_{\rm tot}$ --- but
that is an inference about intent, not a measurement.  The term appears identically in
\code{model401}, \code{model502} and \code{model600} with no comment, so it is worth
confirming with the authors before changing three models.
\end{keybox}

\section{Measurement}

The A/B diagnostic (\code{models/model600/mod\_mach1\_diag.f90}, enabled by
\code{diag\_mach1}) accumulates both forms at every boundary node and reports their ratio.

\begin{codeline}{mod\_boundary\_conditions.f90:632--650 \quad\normalfont\itshape(diagnostic only; imposes nothing)}
\begin{lstlisting}
vE_n_d  =   vE_R_d * normal(1) + vE_Z_d * normal(2)      ! + = into the wall
d_jrk_d = factor * BigR**2 * u0_b / ps0_b
call mach1_diag_add(BigR, factor*cs0, d_jrk_d, vE_n_d, vE_t_d, bn, dl)
\end{lstlisting}
\end{codeline}

\noindent with, inside \code{mach1\_diag\_add}, $\code{d\_slp} = \code{vE\_n}/|\code{bn}|$.
By \eqref{eq:identity} the $u$-dependence cancels identically and
\begin{equation}
  \text{ratio} \;=\; \frac{\code{d\_jrk}}{\code{d\_slp}}
  \;=\; -\operatorname{sign}(b_n)\,\frac{\mathcal{F}}{B_{\rm tot}} ,
  \label{eq:predratio}
\end{equation}
a purely \emph{geometric} quantity.  This is a strong prediction: the ratio should be frozen
in time even while both drift terms evolve.  It is.  The measured min, max and mean were
byte-identical at every output step over the run.

\begin{center}
\begin{tabular}{lcccc}
\toprule
target & $R$ [m] & measured ratio & $R/F_0$ & agreement \\
\midrule
INNER  & $\approx 1.256$ & $+0.431$ & $0.423$ & $2\%$ \\
OUTER  & $\approx 1.594$ & $-0.533$ & $0.536$ & $0.5\%$ \\
\bottomrule
\end{tabular}
\end{center}

\noindent Since $B_{\rm tot} = \sqrt{F_0^2+|\nabla\psi|^2}/R \ge F_0/R$, the quantity
$1/B_{\rm tot}$ must sit just \emph{below} $R/F_0$ --- and the outer measurement does exactly
that.  The opposite signs on the two targets are $\operatorname{sign}(b_n)$ flipping between an
entering and an exiting field line, as \eqref{eq:predratio} requires.  With $F_0 = 2.9723$,
\eqref{eq:ratio} is confirmed to three digits on two independent targets.

\subsection*{Does the difference matter?}

Measured at steady state (weak-form residual $9.8\times10^{-4}$, electrical closure $0.00\%$ on
both targets), in units of $c_s$:

\begin{center}
\begin{tabular}{lcccc}
\toprule
& \multicolumn{2}{c}{JOREK, eq.~\eqref{eq:jorekvpar}} & \multicolumn{2}{c}{required, eq.~\eqref{eq:required}} \\
\cmidrule(lr){2-3}\cmidrule(lr){4-5}
target & mean & max & mean & max \\
\midrule
INNER & $1.36\%$ & $115\%$ & $3.22\%$ & $272\%$ \\
OUTER & $0.25\%$ & $ 46\%$ & $0.48\%$ & $ 85\%$ \\
\bottomrule
\end{tabular}
\end{center}

\noindent A shortfall of $1.9\%$ and $0.2\%$ of $c_s$ \emph{in the mean} --- but up to $150\%$
of $c_s$ \emph{locally}.  The amplifier is grazing incidence:
\begin{equation}
  \langle|b_n|\rangle = 0.0080 \;\;(0.46^\circ), \qquad \min|b_n| = 0.0000 \;\;\text{(tangency)},
\end{equation}
and the correction goes as $\vE\cdot\nhat/b_n$.  It is resolvable
($|\vE\cdot\nhat/b_n| > 10^{-3}c_s$) on only $24\%$ of the inner and $2.3\%$ of the outer
boundary weight: concentrated near tangency, not spread over the target.  The fraction of
boundary weight on which the correction would \emph{reverse} the sign of $v_\parallel$ is
$0.89\%$, so SOLPS's $\pm 2c_s|b_x|$ clip is not optional.

\section{If it is to be changed}
\label{sec:fix}

Equation~\eqref{eq:required} needs the drift term of \eqref{eq:mach1} multiplied by
$B_{\rm tot}$, i.e.\ the $1/B_{\rm tot}$ dropped from that term only:
\begin{equation}
  \frac{\mathcal{F}}{B_{\rm tot}}R^2\frac{\partial_b u}{\partial_b\psi}
  \;\longrightarrow\;
  \mathcal{F}\,R^2\frac{\partial_b u}{\partial_b\psi},
  \qquad\text{clipped to } |\,\cdot\,| \le 2c_s .
\end{equation}
The Jacobian entry \code{Mach1BC\_u} (line 620) scales identically, and
\code{dMach1BC\_ubb} (line 656) likewise --- no new columns are needed, because the derivative
of the term with respect to $u$ already exists.  Put behind a flag
(\code{bohm\_drift\_compatible}, default \code{.false.}) the A/B becomes a namelist switch on a
single build, and the diagnostic provides its own check: ratio \eqref{eq:predratio} must go to
$-\operatorname{sign}(b_n)$ when the flag is on.

\vspace{1em}
\hrule
\vspace{0.6em}
\noindent\footnotesize
Source: \code{models/model600/mod\_boundary\_conditions.f90} and
\code{models/model600/mod\_mach1\_diag.f90}, branch \code{thermoelectric-ohm},
commit \code{9508effe0}.  Run: AUG \#38773, \code{model600}, boundary type 1, $n_{\rm tor}=1$,
weak Galerkin sheath $j$--$V$ BC active.

\end{document}
